
\clearpage
\phantomsection
\addcontentsline{toc}{chapter}{MỞ ĐẦU}
{\LARGE\bfseries MỞ ĐẦU}\par\vspace{1em}

\phantomsection
\addcontentsline{toc}{section}{1. Lý do chọn đề tài}
{\Large\bfseries 1. Lý do chọn đề tài}\par\vspace{0.7em}

FamilyConnect là một nền tảng cộng đồng gia đình số tích hợp AI, cho phép các dòng họ số hóa gia phả, kết nối thành viên, tổ chức sự kiện và lưu trữ di sản gia đình qua nhiều thế hệ. Đây là một bài toán CSDL có độ phức tạp phù hợp để vận dụng đầy đủ quy trình thiết kế đã học trong học phần: từ phân tích yêu cầu, xây dựng mô hình quan niệm ER/EER, chuyển đổi sang mô hình quan hệ, phân tích phụ thuộc hàm và chuẩn hóa, cho đến thiết kế vật lý trên hệ quản trị CSDL cụ thể.

\phantomsection
\addcontentsline{toc}{section}{2. Mục tiêu đồ án}
{\Large\bfseries 2. Mục tiêu đồ án}\par\vspace{0.7em}

\begin{itemize}
\item Vận dụng đúng chu kỳ sống của một CSDL để tổ chức quy trình thiết kế.
\item Xây dựng mô hình thực thể – kết hợp (ERD/EER) đầy đủ theo phương pháp luận: xác định thực thể, thuộc tính, mối kết hợp, bản số, thực thể yếu, chuyên biệt hóa/tổng quát hóa.
\item Phân tích phụ thuộc hàm, chứng minh khóa, chuẩn hóa lược đồ đạt tối thiểu 3NF/BCNF.
\item Thiết kế vật lý: viết script T-SQL triển khai trên Microsoft SQL Server (đúng theo môi trường thực hành Docker ).
\end{itemize}

\phantomsection
\addcontentsline{toc}{section}{3. Phạm vi đồ án}
{\Large\bfseries 3. Phạm vi đồ án}\par\vspace{0.7em}

Đồ án tập trung vào 9 nhóm chức năng của FamilyConnect: Người dùng \& Bảo mật, Gia đình \& Gia phả, Cộng đồng, Sự kiện, Danh bạ gia đình, Di sản gia đình, Dịch vụ AI, Thống kê – Báo cáo, và Quản trị hệ thống. Phần trình bày chi tiết phương pháp (ERD Chen, phân tích FD, chứng minh dạng chuẩn) được thực hiện đầy đủ trên nhóm lõi Gia phả – Cộng đồng – Sự kiện; các nhóm còn lại được thiết kế theo cùng phương pháp luận và trình bày ở mức lược đồ quan hệ hoàn chỉnh kèm khẳng định dạng chuẩn đạt được.